\documentclass[12pt,a4paper]{article}

% ─── Packages ────────────────────────────────────────────────────────────────
\usepackage[utf8]{inputenc}
\usepackage[T1]{fontenc}
\usepackage[french]{babel}
\usepackage{geometry}
\geometry{margin=2.5cm}

\usepackage{amsmath, amssymb, amsthm}
\usepackage{algorithm}
\usepackage{algpseudocode}
\usepackage{tikz}
\usetikzlibrary{arrows.meta, positioning, shapes.geometric,
                decorations.markings, fit, backgrounds, patterns}
\usepackage{pgfplots}
\pgfplotsset{compat=1.18}
\usepackage{booktabs}
\usepackage{array}
\usepackage{xcolor}
\usepackage{hyperref}
\usepackage{enumitem}
\usepackage{mdframed}
\usepackage{lmodern}
\usepackage{microtype}
\usepackage{caption}
\usepackage{tcolorbox}
\tcbuselibrary{skins, breakable}

% ─── Colors ──────────────────────────────────────────────────────────────────
\definecolor{urgence}{RGB}{180,30,30}
\definecolor{theorie}{RGB}{30,80,160}
\definecolor{pratique}{RGB}{20,130,80}
\definecolor{attention}{RGB}{200,120,0}
\definecolor{bleuléger}{RGB}{230,240,255}
\definecolor{vertléger}{RGB}{230,248,238}
\definecolor{twist}{RGB}{20,120,160}
\definecolor{twistléger}{RGB}{220,240,250}
\definecolor{stormy}{RGB}{90,90,130}
\definecolor{stormyléger}{RGB}{235,235,245}

% ─── Custom environments ──────────────────────────────────────────────────────
\newtcolorbox{defbox}[1]{
  colback=bleuléger, colframe=theorie, fonttitle=\bfseries,
  title={#1}, breakable, sharp corners=south
}

\newtcolorbox{metierbox}[1]{
  colback=vertléger, colframe=pratique, fonttitle=\bfseries,
  title={#1}, breakable, sharp corners=south
}

\newtcolorbox{alertbox}[1]{
  colback=yellow!10, colframe=attention, fonttitle=\bfseries,
  title={#1}, breakable, sharp corners=south
}

\newtcolorbox{jurybox}[1]{
  colback=red!5, colframe=urgence, fonttitle=\bfseries,
  title={#1}, breakable, sharp corners=south
}

\newtcolorbox{twistbox}[1]{
  colback=twistléger, colframe=twist, fonttitle=\bfseries\large,
  title={#1}, breakable, sharp corners=south
}

\newtcolorbox{stormbox}[1]{
  colback=stormyléger, colframe=stormy, fonttitle=\bfseries,
  title={#1}, breakable, sharp corners=south
}

% ─── Theorem styles ───────────────────────────────────────────────────────────
\theoremstyle{definition}
\newtheorem{definition}{Définition}[section]
\newtheorem{propriete}{Propriété}[section]
\newtheorem{invariant}{Invariant métier}[section]

\theoremstyle{remark}
\newtheorem{remarque}{Remarque}[section]
\newtheorem{exemple}{Exemple}[section]

% ─── Hyperref setup ───────────────────────────────────────────────────────────
\hypersetup{
  colorlinks=true,
  linkcolor=theorie,
  urlcolor=urgence,
  pdftitle={Twist 6 — Temps de trajet incertain},
  pdfauthor={Système de routage critique},
}

% ─────────────────────────────────────────────────────────────────────────────
\begin{document}

% ═══════════════════════════════════════════════════════════════════
%  COUVERTURE
% ═══════════════════════════════════════════════════════════════════
\begin{titlepage}
  \centering
  \vspace*{2cm}
  {\color{twist}\rule{\linewidth}{2pt}}\\[0.5em]
  {\Huge\bfseries Twist 6}\\[0.4em]
  {\Large\bfseries Le temps de trajet devient incertain\\[0.2em]
  selon l'heure et la météo}\\[0.4em]
  {\color{twist}\rule{\linewidth}{2pt}}\\[2em]

  % Schéma illustratif : même arc, deux distributions selon heure/météo
  \begin{tikzpicture}
    \begin{axis}[
      width=12cm, height=5.5cm,
      xlabel={Temps de trajet estimé (minutes)},
      ylabel={Densité de probabilité},
      xmin=4, xmax=28,
      ymin=0, ymax=0.26,
      ytick=\empty,
      grid=major, grid style={gray!20},
      legend pos=north east,
      legend style={font=\small},
      title={\small Distribution du temps de trajet — même arc, mêmes conditions nominales}
    ]
    % Heure creuse, temps sec — distribution étroite
    \addplot[color=pratique, thick, smooth, domain=4:28, samples=80]
      {0.22*exp(-0.5*((x-8)/1.5)^2)};
    \addlegendentry{07h00, sec \quad $\bar{t}=8$ min, $\sigma=1.5$}

    % Heure de pointe, pluie — distribution large, décalée
    \addplot[color=urgence, thick, smooth, dashed, domain=4:28, samples=80]
      {0.1*exp(-0.5*((x-16)/3.8)^2)};
    \addlegendentry{08h00, pluie \quad $\bar{t}=16$ min, $\sigma=3.8$}

    % Percentile 90 pour pluie
    \addplot[color=attention, thick, dotted, domain=21.87:21.87]
      {0} -- (axis cs:21.87, 0.22);
    \node[font=\footnotesize, color=attention, rotate=90]
      at (axis cs:22.5, 0.11) {$P_{90} = 21$ min};
    \end{axis}
  \end{tikzpicture}\\[1.5em]

  {\large Note technique — Routage robuste sous incertitude stochastique}\\[1em]
  {\large\color{gray} Extension du moteur de routage dynamique A*}\\[2.5em]

  \begin{tabular}{rl}
    \textbf{Sujet :}       & Sujet 07 — Hackverse 2026 \\
    \textbf{Twist :}       & \#6 — Temps de trajet incertain \\
    \textbf{Sources d'incertitude :} &
      Heure de la journée \textbf{+} état météorologique \\
    \textbf{Rupture clé :} &
      Les poids deviennent des variables aléatoires, pas des scalaires \\
  \end{tabular}

  \vfill
  {\small\color{gray} Document pédagogique — Présentation jury et équipe technique}
\end{titlepage}

% ═══════════════════════════════════════════════════════════════════
\tableofcontents
\newpage

% ═══════════════════════════════════════════════════════════════════
\section*{Introduction}
\addcontentsline{toc}{section}{Introduction}

Dans le moteur présenté dans \texttt{routage\_dynamique.tex}, le poids
d'un arc est une fonction déterministe du temps :
\[
  w(u,v,t) = w_0(u,v) \cdot \mu(u,v,t) \cdot \rho(u,v,t) \cdot \phi_\tau(u,v)
\]
À l'instant $t$ donné, le système \emph{connaît} $\mu$ et $\rho$.
Il choisit donc le chemin qui minimise une somme de scalaires certains.

\begin{twistbox}{Twist 6 — Le problème change de nature mathématique}
Le temps de trajet réel sur un arc \textbf{n'est pas un scalaire connu}
à l'avance. C'est une \textbf{variable aléatoire} dont la distribution
dépend de l'heure et de la météo.

\medskip
\begin{itemize}[leftmargin=1.5em]
  \item À \textbf{8h00 sous la pluie}, l'arc Kennedy prend entre 10 et
    22 minutes selon les imprévus : feux en panne, marché débordant, poids
    lourds glissants.
  \item À \textbf{3h00 par temps sec}, ce même arc prend entre 6 et
    9 minutes — distribution étroite et prévisible.
\end{itemize}

\medskip
Le moteur ne peut plus simplement calculer « quel chemin minimise la
somme des temps ». Il doit répondre à une question plus difficile :
\textbf{quel chemin offre la meilleure garantie d'arrivée dans les délais,
compte tenu de ce que l'on ne sait pas ?}
\end{twistbox}

Ce twist transforme le problème de routage en un problème
d'\textbf{optimisation stochastique} : choisir sous incertitude, en
définissant explicitement ce que l'on cherche à garantir.

\newpage

% ═══════════════════════════════════════════════════════════════════
\section{Ce qui change par rapport au modèle déterministe}

\subsection{Dans le modèle déterministe (twists précédents)}

Le multiplicateur $\mu(u,v,t)$ était un \textbf{scalaire unique} issu
de la table de trafic, mis à jour toutes les 5 minutes. Le système
savait : « à 08h05, cet arc prend 14 minutes ». L'incertitude était
absente de la formule.

\subsection{Dans le modèle stochastique (twist 6)}

Le multiplicateur $\mu$ devient une \textbf{variable aléatoire}
$M(u,v,t)$. La table de trafic ne fournit plus un seul chiffre, mais
une distribution : moyenne $\bar{\mu}$, écart-type $\sigma_\mu$,
forme de la courbe. L'incertitude est explicitement encodée.

\begin{center}
\begin{tabular}{lcc}
  \toprule
  & \textbf{Modèle déterministe} & \textbf{Modèle stochastique} \\
  \midrule
  Poids d'un arc       & scalaire $w(u,v,t) \in \mathbb{R}^+$
                       & variable aléatoire $W(u,v,t,\theta)$ \\
  Coût d'un chemin     & somme de scalaires
                       & somme de variables aléatoires \\
  Critère d'optimalité & minimiser la somme
                       & minimiser une mesure de risque \\
  Outil algorithmique  & A* classique
                       & A* robuste (§\ref{sec:astar_robuste}) \\
  \bottomrule
\end{tabular}
\end{center}

\subsection{Deux sources d'incertitude distinctes}

\begin{defbox}{Source 1 — Heure de la journée}
L'incertitude temporelle traduit la \textbf{variabilité intra-heure} :
même en sachant que l'on est à 08h10 un lundi, le trafic sur un arc
donné peut prendre des valeurs très différentes selon les jours,
les micro-événements et la réaction des conducteurs. L'écart-type
$\sigma_\mu$ est élevé aux heures de pointe (7h–9h, 17h–20h) et
faible la nuit (22h–5h).
\end{defbox}

\begin{defbox}{Source 2 — État météorologique}
L'incertitude météorologique a deux composantes :
\begin{enumerate}[leftmargin=1.5em]
  \item \textbf{L'incertitude sur l'état courant} : le capteur météo
    dit « pluie légère » mais ne quantifie pas l'intensité exacte ni
    l'état exact de la chaussée.
  \item \textbf{L'incertitude sur l'état futur} : une prévision à
    30 minutes donne « 70\% de risque de pluie forte ». Pendant le
    trajet, le temps peut évoluer. Un arc sec au départ peut être
    mouillé à l'arrivée.
\end{enumerate}
\end{defbox}

\newpage

% ═══════════════════════════════════════════════════════════════════
\section{Modélisation des poids stochastiques}

\subsection{Poids d'un arc comme variable aléatoire}

\begin{defbox}{Définition : Poids stochastique d'un arc}
Le temps de trajet sur l'arc $(u,v)$ à l'instant $t$ sous conditions
météorologiques $\theta$ est modélisé comme :
\[
  W(u,v,t,\theta) = w_0(u,v) \cdot M(u,v,t) \cdot R(u,v,\theta)
\]
où :
\begin{itemize}[leftmargin=1.5em]
  \item $w_0(u,v)$ est le \textbf{temps de base} (conditions idéales,
    scalaire constant) ;
  \item $M(u,v,t)$ est le \textbf{multiplicateur temporel stochastique},
    variable aléatoire dont les paramètres dépendent de $t$ ;
  \item $R(u,v,\theta)$ est le \textbf{facteur météorologique stochastique},
    variable aléatoire dépendant de l'état météo $\theta$.
\end{itemize}
\end{defbox}

\subsection{Distribution du multiplicateur temporel $M(u,v,t)$}

Sur la base des 2\,400 relevés du fichier
\texttt{traffic\_by\_segment\_hour.csv}, le multiplicateur $M$
est modélisé par une \textbf{loi log-normale}, naturelle pour
les temps de trajet (toujours positifs, asymétrie vers les
valeurs élevées) :
\[
  \ln M(u,v,t) \sim \mathcal{N}\!\left(\ln\bar{\mu}(u,v,t),\;
  \sigma_\mu^2(u,v,t)\right)
\]

Les paramètres $\bar{\mu}$ (moyenne) et $\sigma_\mu$ (écart-type)
varient avec l'heure :

\begin{center}
\begin{tabular}{lcccc}
  \toprule
  \textbf{Plage horaire} & \textbf{$\bar{\mu}$ (axe primaire)} &
  \textbf{$\sigma_\mu$} & \textbf{$P_{90}$} & \textbf{Interprétation} \\
  \midrule
  Nuit (22h–5h)         & $\times 1{,}2$ & $0{,}08$ & $\times 1{,}3$ &
    Trafic rare, imprévisibilité faible \\
  Heure creuse (9h–16h) & $\times 1{,}5$ & $0{,}18$ & $\times 1{,}8$ &
    Trafic modéré, variance modérée \\
  Pointe mat. (7h–9h)   & $\times 2{,}1$ & $0{,}35$ & $\times 3{,}0$ &
    Congestion forte, très imprévisible \\
  Pointe soir (17h–20h) & $\times 2{,}3$ & $0{,}40$ & $\times 3{,}3$ &
    Pire plage horaire, variance maximale \\
  \bottomrule
\end{tabular}
\end{center}

\begin{alertbox}{Lecture de la colonne $P_{90}$}
Le $P_{90}$ est le \textbf{90\textsuperscript{e} percentile} :
dans 90\% des cas, le multiplicateur réel sera \emph{inférieur} à
cette valeur. En pointe du soir, $P_{90} = \times 3{,}3$ signifie
qu'une ambulance doit planifier son trajet comme si les conditions
étaient 3,3 fois plus lentes que l'idéal — et s'attendre à ce que
1 fois sur 10, elles soient encore pires.
\end{alertbox}

\subsection{Distribution du facteur météorologique $R(u,v,\theta)$}

L'état météorologique $\theta$ est discrétisé en trois niveaux,
chacun associé à une distribution sur le facteur de ralentissement :

\begin{center}
\begin{tabular}{lccc}
  \toprule
  \textbf{État $\theta$} & \textbf{$\mathbb{E}[R]$} &
  \textbf{$\sigma_R$} & \textbf{Impact principal} \\
  \midrule
  Sec ($\theta_0$)          & $1{,}00$ & $0{,}05$ &
    Trafic nominal, variance faible \\
  Pluie légère ($\theta_1$) & $1{,}25$ & $0{,}15$ &
    Ralentissement modéré, variabilité accrue \\
  Pluie forte ($\theta_2$)  & $1{,}60$ & $0{,}30$ &
    Aquaplanage, visibilité réduite, accidents \\
  \bottomrule
\end{tabular}
\end{center}

\noindent L'état futur $\theta$ est connu avec une probabilité $p(\theta)$
fournie par la prévision météo. Par exemple :
\[
  p(\theta_0) = 0{,}15, \quad
  p(\theta_1) = 0{,}55, \quad
  p(\theta_2) = 0{,}30
\]
\og{}15\% de chance que ça reste sec, 55\% de pluie légère,
30\% de pluie forte dans les 30 prochaines minutes.\fg{}

\subsection{Distribution du coût total d'un chemin}

Pour un chemin $P = (v_0, v_1, \ldots, v_k)$, le temps de trajet
total est la somme des poids stochastiques des arcs :
\[
  T(P) = \sum_{i=0}^{k-1} W(v_i, v_{i+1},\, t_i,\, \theta)
\]

Par le \textbf{théorème de la limite centrale} (applicable dès
$k \geq 5$ arcs), $T(P)$ est approximativement log-normale,
avec :
\[
  \mathbb{E}[T(P)] = \sum_{i} \mathbb{E}[W_i], \qquad
  \text{Var}[T(P)] \approx \sum_{i} \text{Var}[W_i]
  \quad\text{(si arcs peu corrélés)}
\]

\begin{alertbox}{Corrélation météo entre arcs voisins}
Les arcs géographiquement proches sont affectés par la \emph{même}
pluie au même moment : leurs facteurs $R$ sont corrélés. Ignorer
cette corrélation sous-estime la variance totale. En pratique,
on la traite par \textbf{clustering géographique} : les arcs d'un
même quartier partagent le même tirage de $\theta$, ceux de quartiers
différents tirent indépendamment.
\end{alertbox}

\newpage

% ═══════════════════════════════════════════════════════════════════
\section{Critères de décision sous incertitude}
\label{sec:criteres}

\subsection{Le problème du critère}

Minimiser $\mathbb{E}[T(P)]$ revient à ignorer le risque : deux chemins
avec la même espérance mais des variances très différentes sont traités
à l'égalité. Pour un véhicule de secours, cette équivalence est inacceptable.

\begin{center}
\begin{tikzpicture}
  \begin{axis}[
    width=12cm, height=5cm,
    xlabel={Temps de trajet (min)},
    ylabel={Densité},
    xmin=5, xmax=30,
    ymin=0, ymax=0.32,
    ytick=\empty,
    grid=major, grid style={gray!20},
    legend pos=north east, legend style={font=\small},
    title={\small Deux chemins de même espérance (12 min) — variances très différentes}
  ]
  % Chemin A — variance faible (sûr)
  \addplot[color=pratique, thick, smooth, domain=5:30, samples=80]
    {0.30*exp(-0.5*((x-12)/1.8)^2)};
  \addlegendentry{Chemin A — $\sigma=1{,}8$ min (fiable)}

  % Chemin B — variance forte (risqué)
  \addplot[color=urgence, thick, smooth, dashed, domain=5:30, samples=80]
    {0.12*exp(-0.5*((x-12)/4.5)^2)};
  \addlegendentry{Chemin B — $\sigma=4{,}5$ min (risqué)}

  % Ligne médiane
  \addplot[color=gray, thin, dotted] coordinates {(12,0)(12,0.31)};
  \node[font=\footnotesize, color=gray] at (axis cs:13.2,0.28)
    {$\mathbb{E}=12$};
  \end{axis}
\end{tikzpicture}
\end{center}

Si la fenêtre d'intervention critique est de 15 minutes, le chemin A a
une probabilité d'arrivée à temps de $\approx$\,95\%, contre
$\approx$\,75\% pour le chemin B — malgré la même espérance.

\subsection{Critère retenu : minimisation au percentile $P_{90}$}

\begin{defbox}{Définition : Critère du percentile $P_\alpha$}
Le chemin optimal est celui qui minimise le \textbf{$P_\alpha$-quantile}
du temps de trajet :
\[
  P^* = \arg\min_{P \in \mathcal{P}(s,h)} \;
  \text{VaR}_\alpha[T(P)]
\]
où $\text{VaR}_\alpha[T(P)] = \inf\{t \mid \mathbb{P}[T(P) \leq t] \geq \alpha\}$.

En pratique, avec $\alpha = 0{,}90$ : on garantit que dans 90\%
des conditions réalisées, l'ambulance arrivera en moins de
$\text{VaR}_{0.90}[T(P^*)]$ minutes.
\end{defbox}

Pour une loi log-normale de paramètres $(\mu_P, \sigma_P)$ :
\[
  \text{VaR}_{0.90}[T(P)] = e^{\mu_P + 1{,}282\,\sigma_P}
\]

Ce critère est réglable selon l'urgence du cas : $\alpha = 0{,}80$
pour une intervention non critique, $\alpha = 0{,}95$ pour un arrêt
cardiaque.

\subsection{Critère alternatif : pénalisation de la variance}

\begin{defbox}{Définition : Critère risque-espérance (mean-variance)}
\[
  P^* = \arg\min_{P} \Bigl[
    \mathbb{E}[T(P)] \;+\; \lambda \cdot \sigma[T(P)]
  \Bigr]
\]
Le paramètre $\lambda \geq 0$ contrôle l'aversion au risque.
$\lambda = 0$ revient au critère espérance pure ; $\lambda = 2$
donne un compromis standard (règle des deux sigma).
\end{defbox}

Pour les véhicules d'urgence, le critère $P_{90}$ est préféré car
il a une signification opérationnelle claire (« dans 9 cas sur 10,
on arrive dans les délais »), là où $\lambda$ est un paramètre
abstrait difficile à calibrer et à expliquer.

\newpage

% ═══════════════════════════════════════════════════════════════════
\section{Algorithme A* robuste}
\label{sec:astar_robuste}

\subsection{Principe : transformer les poids stochastiques en poids scalaires conservatifs}

L'idée centrale est de remplacer chaque poids stochastique par un
\textbf{poids conservatif scalaire} qui encode à la fois l'espérance
et le risque. Ainsi, A* reste applicable sans modification de sa
structure de file de priorité.

\begin{defbox}{Définition : Poids conservatif $\tilde{w}$}
Le poids conservatif d'un arc $(u,v)$ au niveau de confiance $\alpha$
est défini comme le $\alpha$-quantile de sa distribution :
\[
  \tilde{w}_\alpha(u,v,t,\theta) =
  \mathbb{E}[W(u,v,t,\theta)] \;+\; z_\alpha \cdot \sigma[W(u,v,t,\theta)]
\]
où $z_\alpha$ est le quantile de la loi normale standard :
$z_{0.80} = 0{,}84$, $z_{0.90} = 1{,}28$, $z_{0.95} = 1{,}64$.
\end{defbox}

Ce remplacement est une \textbf{approximation conservative} : la somme
des quantiles individuels majore le quantile de la somme. Le chemin
obtenu n'est pas le strict optimal $P_{90}$, mais il en est une borne
supérieure garantie — propriété précieuse pour un système critique.

\subsection{Pseudocode de A* robuste}

\begin{algorithm}[H]
\caption{A* Robuste — poids stochastiques au niveau $\alpha$}
\begin{algorithmic}[1]
\Require Graphe $G$, source $s$, destination $h^*$, instant $t_0$,
         prévision météo $\mathbf{p}(\theta)$, niveau $\alpha$
\Ensure Chemin robuste $P^*$ et son $P_\alpha$-coût estimé

\State \textbf{Précalculer} pour chaque arc $(u,v) \in E$ :
\State \quad $\bar{W}(u,v) \gets w_0(u,v) \cdot
       \mathbb{E}[M(u,v,t_0)] \cdot \sum_\theta p(\theta)\,\mathbb{E}[R(u,v,\theta)]$
\State \quad $\sigma_W(u,v) \gets w_0(u,v) \cdot
       \sqrt{\text{Var}[M \cdot R](u,v,t_0,\mathbf{p})}$
\State \quad $\tilde{w}(u,v) \gets \bar{W}(u,v) + z_\alpha \cdot \sigma_W(u,v)$
       \Comment{Poids conservatif}

\Statex
\State Lancer A*$(G,\, s,\, h^*,\, \tilde{w})$
       \Comment{A* standard avec les poids $\tilde{w}$}
\State $P^* \gets$ chemin retourné par A*
\State \textbf{journaliser}$\bigl(t_0,\; \alpha,\; \mathbf{p}(\theta),\;
       P^*,\; \tilde{w}[P^*],\; \bar{W}[P^*],\; \sigma_W[P^*]\bigr)$
\State \Return $P^*$
\end{algorithmic}
\end{algorithm}

\subsection{Heuristique admissible en contexte stochastique}

Pour que A* reste optimal sur les poids conservatifs $\tilde{w}$,
l'heuristique $h(n)$ doit majorer le poids conservatif du chemin
restant — pas seulement le temps minimal. La condition d'admissibilité
devient :
\[
  h(n) \leq \tilde{w}^*(n, h^*) \quad \forall n \in V
\]

La borne suivante est admissible et facile à calculer :
\[
  h(n) = \frac{d_{\text{euclidienne}}(n, h^*)}{v_{\max}}
  \cdot (1 + z_\alpha \cdot \sigma_{\min})
\]
où $\sigma_{\min}$ est l'écart-type minimal observé sur les arcs du réseau.
Le facteur $(1 + z_\alpha \cdot \sigma_{\min})$ garantit que l'heuristique
intègre une part d'incertitude sans surestimer.

\subsection{Recalcul dynamique des poids}

Les poids conservatifs $\tilde{w}$ doivent être recalculés lors de
deux types d'événements :

\begin{enumerate}[leftmargin=2em, itemsep=4pt]
  \item \textbf{Toutes les 5 minutes} — recalcul planifié, comme dans
    le modèle déterministe de base.
  \item \textbf{Sur mise à jour météo significative} — si la prévision
    $\mathbf{p}(\theta)$ change d'au moins $\Delta p = 0{,}15$ en
    probabilité (par exemple, passage de 30\% à 60\% de pluie forte),
    recalcul immédiat déclenché.
\end{enumerate}

\begin{metierbox}{Pourquoi le seuil $\Delta p = 0{,}15$ ?}
Un changement de prévision inférieur à 15 points de probabilité
modifie les poids conservatifs de moins de 4\% en moyenne —
insuffisant pour changer le chemin optimal dans la grande majorité
des cas. Ce seuil évite des recalculs incessants pour des variations
météo mineures, tout en réagissant aux évolutions significatives.
\end{metierbox}

\newpage

% ═══════════════════════════════════════════════════════════════════
\section{Scénario : l'ambulance de Bépanda — version météo}

Ce scénario reprend le cas de référence et y intègre l'incertitude
temporelle et météorologique.

\subsection{Contexte — appel à 07h58, lundi matin}

Il est 07h58 un lundi. Un enfant de sept ans fait des convulsions à
Bépanda. L'appel arrive au SAMU. Le système reçoit simultanément :

\begin{itemize}[leftmargin=1.5em]
  \item \textbf{Heure} : 07h58 — début de la pointe matinale.
    $\bar{\mu} = 2{,}1$, $\sigma_\mu = 0{,}35$.
  \item \textbf{Prévision météo} : $p(\theta_0) = 0{,}10$,
    $p(\theta_1) = 0{,}60$, $p(\theta_2) = 0{,}30$.
    Pluie légère à forte probable dans 15 minutes.
\end{itemize}

\noindent
Le système calcule les poids conservatifs ($\alpha = 0{,}90$) pour
les deux itinéraires candidats vers l'Hôpital Central :

\medskip
\begin{center}
\begin{tabular}{lccccc}
  \toprule
  \textbf{Itinéraire} & \textbf{$\bar{W}$} &
  \textbf{$\sigma_W$} & \textbf{$\tilde{w}_{90}$} &
  \textbf{Retenu ?} \\
  \midrule
  Bépanda $\to$ Bld Liberté $\to$ Central &
    12,4 min & 3,8 min & 17,3 min & \textbf{oui} \\
  Bépanda $\to$ Makepe $\to$ Ngousso $\to$ Central &
    15,1 min & 1,9 min & 17,5 min & non \\
  \bottomrule
\end{tabular}
\end{center}

\medskip
\noindent
Le premier itinéraire (Boulevard de la Liberté) a une espérance
plus faible (12,4 vs 15,1 min) mais une variance plus élevée
(3,8 vs 1,9 min). Au niveau $P_{90}$, les deux sont presque
équivalents (17,3 vs 17,5 min). Le système retient le Boulevard
de la Liberté pour son espérance meilleure, la différence de
risque étant négligeable.

\begin{stormbox}{Ce que ce calcul révèle}
Sans modèle stochastique, un système naïf aurait simplement
pris le chemin le plus court en espérance (12,4 min) sans
évaluer le risque. Ce n'est que le $P_{90}$ qui révèle que
l'itinéraire « plus long » via Makepe est en réalité presque
aussi bon — et parfois préférable si la tolérance au risque
change.
\end{stormbox}

\subsection{Dégradation météo en cours de trajet — 08h04}

L'ambulance roule depuis 6 minutes sur le Boulevard de la Liberté.
À 08h04, la prévision météo est mise à jour :

\[
  p(\theta_0) = 0{,}05, \quad
  p(\theta_1) = 0{,}25, \quad
  p(\theta_2) = \mathbf{0{,}70}
\]

La probabilité de pluie forte passe de 30\% à 70\% — variation de
$+0{,}40 > \Delta p = 0{,}15$. Le recalcul est déclenché
immédiatement.

\medskip
\begin{center}
\begin{tabular}{lccccc}
  \toprule
  \textbf{Itinéraire (depuis position courante)} &
  \textbf{$\bar{W}$} & \textbf{$\sigma_W$} &
  \textbf{$\tilde{w}_{90}$} & \textbf{Retenu ?} \\
  \midrule
  Continuer Bld Liberté $\to$ Central &
    9,0 min & 4,2 min & 14,4 min & non \\
  Dévier $\to$ Rue Makepe $\to$ Central &
    11,5 min & 2,1 min & \textbf{14,2 min} & \textbf{oui} \\
  \bottomrule
\end{tabular}
\end{center}

\medskip
\noindent
Le Boulevard de la Liberté est une route goudronnée large :
sous pluie forte, sa variance explose (aquaplanage, accidents en
chaîne). La Rue Makepe, piste en terre, est certes plus longue en
espérance — mais sous pluie forte, son trafic faible la rend
\textbf{plus prévisible}. Au $P_{90}$, elle est maintenant
légèrement meilleure (14,2 vs 14,4 min).

\medskip
Le gain est de 0,2 minute, inférieur au seuil de stabilité
$\Delta_{\min} = 2$ minutes. \textbf{Le système maintient le Boulevard
de la Liberté}, note la re-vérification dans le journal, et signale
au conducteur que les conditions se dégradent.

\begin{metierbox}{Décision à 08h04}
Le système ne change pas le chemin mais envoie une alerte au
conducteur : « Pluie forte probable — adaptez votre vitesse.
Itinéraire de repli disponible via Makepe si conditions trop
dégradées. » L'humain reste en mesure de décider si la déviation
vaut l'écart de 2,3 minutes d'espérance supplémentaire.
\end{metierbox}

\subsection{Résolution — 08h15}

L'ambulance arrive à l'Hôpital Central à 08h15. Soit 17 minutes
après le départ — dans la fenêtre $P_{90}$ estimée à 17,3 min.
Le modèle était juste : la pluie a bien ralenti le Boulevard, mais
moins que le scénario catastrophe qui justifiait la déviation.

\newpage

% ═══════════════════════════════════════════════════════════════════
\section{Invariants métier et cas limites}

\subsection{Invariants introduits par le twist 6}

\begin{invariant}[Niveau de confiance minimum]
Le système ne route jamais un véhicule d'urgence sur un itinéraire
dont le $P_{90}$ dépasse le délai critique $T_{\max}$ défini par
le protocole médical du cas, sauf si aucun itinéraire dans le
réseau ne respecte cette contrainte.
\end{invariant}

\begin{invariant}[Recalcul sur événement météo significatif]
Toute mise à jour de prévision météo modifiant $p(\theta)$ de plus
de $\Delta p = 0{,}15$ pour au moins un état $\theta$ déclenche un
recalcul des poids conservatifs et, si nécessaire, de l'itinéraire.
\end{invariant}

\begin{invariant}[Transparence du niveau de confiance]
Chaque itinéraire transmis au conducteur est accompagné de son
$P_{90}$ estimé et de son espérance $\mathbb{E}[T]$, afin que
l'opérateur comprenne la fiabilité de l'estimation.
\end{invariant}

\subsection{Cas limite : variance maximale — le chemin le plus court est dangereux}

Si tous les itinéraires possibles ont un $P_{90}$ supérieur à
$T_{\max}$, le système applique la stratégie de repli suivante :

\begin{enumerate}[leftmargin=2em, itemsep=4pt]
  \item \textbf{Retenir le chemin minimisant l'espérance} (et non
    le $P_{90}$) — en situation sans issue garantie, minimiser le
    temps attendu est la seule stratégie rationnelle.
  \item \textbf{Alerter l'opérateur} avec le message :
    « Aucun itinéraire ne garantit l'arrivée dans les délais.
    Meilleure option en espérance : [itinéraire]. Délai estimé $P_{90}$ : [valeur]. »
  \item \textbf{Activer la priorisation de feux} si le système est
    connecté à la gestion du trafic urbain — pour dégrader artificiellement
    moins les routes à fort trafic.
\end{enumerate}

\subsection{Cas limite : prévision météo indisponible}

Si l'API météo est hors service, le système applique la distribution
\textbf{la plus défavorable observée historiquement} à l'heure courante
sur les 30 derniers jours. C'est une stratégie maximin (minimiser le
coût dans le pire cas) — conservatrice mais jamais naïve.

\begin{alertbox}{Signalisation à l'opérateur}
En cas de données météo manquantes, l'interface affiche :
« Données météo indisponibles — calcul sur historique conservatif.
Les temps estimés peuvent être surestimés de 15 à 30\%. »
\end{alertbox}

\newpage

% ═══════════════════════════════════════════════════════════════════
\section{Implémentation dans le moteur existant}

\subsection{Extension du backend Python}

Le twist 6 s'intègre dans la couche de pondération existante
sans modifier l'algorithme A* lui-même.

\begin{center}
\begin{tabular}{>{\ttfamily}p{4.5cm} p{8.5cm}}
  \toprule
  \textbf{Fichier} & \textbf{Modification} \\
  \midrule
  routing/weights.py &
    Nouveau module : calcul de $\bar{W}$, $\sigma_W$ et $\tilde{w}_\alpha$
    pour chaque arc, à partir de la table de trafic et de la prévision
    météo courante. \\[0.4em]
  routing/algorithms.py &
    Remplacer l'appel à \texttt{compute\_route(G, s, h)} par
    \texttt{compute\_robust\_route(G, s, h, alpha=0.90)} qui précalcule
    les $\tilde{w}$ avant de lancer A*. \\[0.4em]
  api/weather\_feed.py &
    Nouveau module de polling de l'API météo (Open-Meteo, toutes les
    60 s). Expose \texttt{get\_weather\_forecast(lat, lon, horizon\_min)}
    retournant $\mathbf{p}(\theta)$. \\[0.4em]
  routing/graph.py &
    Étendre le schéma de chaque arc avec les champs \texttt{mu\_mean},
    \texttt{mu\_sigma} (par plage horaire) et \texttt{r\_mean},
    \texttt{r\_sigma} (par état météo). \\[0.4em]
  routing/logger.py &
    Étendre le log de décision avec les champs de dispersion :
    \texttt{w\_mean}, \texttt{w\_sigma}, \texttt{w\_p90},
    \texttt{alpha\_used}, \texttt{weather\_proba}. \\
  \bottomrule
\end{tabular}
\end{center}

\subsection{Exemple de calcul du poids conservatif en Python}

\begin{verbatim}
def conservative_weight(edge, t, weather_proba, alpha=0.90):
    """
    Calcule le poids conservatif d'un arc au niveau alpha.
    edge     : dict avec w0, mu_mean[heure], mu_sigma[heure],
                         r_mean[état], r_sigma[état]
    t        : heure courante (0-23)
    weather_proba : dict {0: p_sec, 1: p_léger, 2: p_fort}
    """
    from scipy.stats import norm
    import numpy as np

    hour = int(t)
    w0      = edge['w0']
    mu_m    = edge['mu_mean'][hour]
    mu_s    = edge['mu_sigma'][hour]

    # Espérance et variance de R (mélange météo)
    r_m = sum(weather_proba[k] * edge['r_mean'][k] for k in range(3))
    r_v = sum(weather_proba[k] * (edge['r_sigma'][k]**2
              + edge['r_mean'][k]**2) for k in range(3)) - r_m**2

    # Propagation des incertitudes (produit de v.a. indépendantes)
    W_mean  = w0 * mu_m * r_m
    W_sigma = w0 * np.sqrt(mu_s**2 * r_m**2 + r_v * mu_m**2
                           + mu_s**2 * r_v)

    z_alpha = norm.ppf(alpha)   # z_0.90 = 1.282
    return W_mean + z_alpha * W_sigma
\end{verbatim}

\subsection{Extension de la réponse API}

L'endpoint \texttt{/api/route/} enrichit sa réponse :

\begin{verbatim}
{
  "route": { ... },
  "uncertainty": {
    "alpha":            0.90,
    "expected_minutes": 12.4,
    "sigma_minutes":    3.8,
    "p90_minutes":      17.3,
    "weather_forecast": { "dry": 0.10, "light": 0.60, "heavy": 0.30 },
    "weather_source":   "Open-Meteo API",
    "weather_freshness_sec": 45
  }
}
\end{verbatim}

\newpage

% ═══════════════════════════════════════════════════════════════════
\section{Communication au jury}

\begin{jurybox}{Message clé pour le jury}
\textbf{Avec ce twist, le système ne dit plus « vous arriverez en
12 minutes ». Il dit « vous arriverez en moins de 17 minutes dans
90\% des cas, compte tenu de la pluie probable et de l'heure de
pointe ».}

\medskip
C'est la différence entre un GPS classique — qui vous annonce un
temps d'arrivée précis et se trompe régulièrement — et un système
d'aide à la décision pour les urgences, qui quantifie honnêtement
ce qu'il ne sait pas.
\end{jurybox}

\subsection{Questions anticipées du jury}

\paragraph{Question : \og{}Pourquoi ne pas simplement ajouter une marge
fixe de sécurité ?\fg{}}

\noindent\textbf{Réponse :} Une marge fixe (« toujours ajouter 5 minutes »)
ignore que l'incertitude est variable : la nuit, elle est faible et la
marge serait inutilement pessimiste ; en pointe sous pluie forte, elle
serait insuffisante. Le modèle stochastique calcule la bonne marge
pour chaque situation réelle, sans sur-correction ni sous-correction.

\paragraph{Question : \og{}Le modèle est compliqué — pouvez-vous l'expliquer
simplement ?\fg{}}

\noindent\textbf{Réponse :} Le système a deux questions à répondre pour
chaque segment de route. Première question : « En moyenne, combien de
temps ça prend à cette heure ? » Deuxième question : « Combien peut-on
se tromper sous ces conditions météo ? » Il additionne les deux pour
obtenir une estimation prudente. C'est exactement ce qu'un conducteur
expérimenté fait mentalement — le système le fait avec des chiffres.

\paragraph{Question : \og{}Et si la prévision météo elle-même est fausse ?\fg{}}

\noindent\textbf{Réponse :} La prévision est intégrée comme une
distribution de probabilités, pas comme une certitude. Si elle prédit
« 70\% de pluie forte », le modèle intègre aussi les 30\% restants.
La robustesse vient précisément du fait que l'on ne suppose pas que
la prévision est parfaite — on la traite comme une source parmi d'autres,
avec son propre degré d'incertitude. En cas de panne totale de la
prévision, le système bascule automatiquement sur l'historique conservatif.

\subsection{Tableau de bord : log enrichi}

\begin{center}
\begin{tabular}{ll}
  \toprule
  \textbf{Champ du log} & \textbf{Exemple de valeur} \\
  \midrule
  Timestamp de calcul       & 2026-04-26 07:58:12 \\
  Véhicule                  & Ambulance SAMU-07 \\
  Heure de départ           & 07h58 (pointe matinale) \\
  Prévision météo           & sec 10\%, léger 60\%, fort 30\% \\
  Niveau de confiance       & $\alpha = 0{,}90$ \\
  Chemin retenu             & Bépanda $\to$ Bld Liberté $\to$ Central \\
  Espérance $\mathbb{E}[T]$ & 12,4 min \\
  Écart-type $\sigma[T]$    & 3,8 min \\
  $P_{90}$ estimé           & 17,3 min \\
  Chemin alternatif évalué  & via Makepe — $P_{90} = 17{,}5$ min \\
  Recalcul déclenché à      & 08h04 (hausse pluie forte : $+40\%$) \\
  Décision recalcul         & Maintien (gain $<\Delta_{\min}$) \\
  Temps réel d'arrivée      & 17 min — dans la fenêtre $P_{90}$ \\
  \bottomrule
\end{tabular}
\end{center}

\newpage

% ═══════════════════════════════════════════════════════════════════
\section*{Conclusion}
\addcontentsline{toc}{section}{Conclusion}

Le twist 6 impose au moteur de routage de \textbf{raisonner sur ses
propres incertitudes}. Ce n'est plus seulement un calculateur
d'itinéraires — c'est un système qui quantifie ce qu'il sait, ce qu'il
ne sait pas, et construit ses décisions en tenant compte explicitement
des deux.

\medskip
Trois acquis structurants émergent de cette évolution :

\begin{enumerate}[leftmargin=2em, itemsep=4pt]
  \item \textbf{Les poids sont des distributions, pas des scalaires.}
    Chaque arc est décrit par une espérance et un écart-type, fonctions
    de l'heure et de la météo. L'algorithme travaille sur des poids
    conservatifs qui agrègent les deux.

  \item \textbf{Le critère d'optimisation est opérationnel.}
    Minimiser le $P_{90}$ répond directement à la question du médecin :
    « Avec quelle probabilité mon patient arrive-t-il dans les délais ? »
    C'est un langage que les opérateurs et les responsables comprennent.

  \item \textbf{La transparence est totale.}
    Chaque décision est accompagnée de son intervalle d'incertitude.
    Si une ambulance arrive en retard, le journal montre si c'était
    un cas dans les 90\% prévus — ou dans les 10\% que le système
    déclarait lui-même ne pas pouvoir garantir.
\end{enumerate}

\medskip
\begin{center}
  {\color{twist}\rule{0.5\linewidth}{1pt}}\\[0.5em]
  {\itshape \og{}Un système honnête ne dit pas « vous arriverez dans 12 minutes ».\\
  Il dit « dans 9 cas sur 10, vous arriverez dans 17 minutes ».\\
  Et dans le 10\textsuperscript{e} cas, il a déjà préparé le plan de repli.\fg{}}\\[0.5em]
  {\color{twist}\rule{0.5\linewidth}{1pt}}
\end{center}

\end{document}